% ============================================================
% Chapter 3: Threat Actors, Motivations and Attack Lifecycle Models
% Language: Greek — edit freely; regenerate from src/content if preferred.
% ============================================================
\chapter{Δράστες Απειλών, Κίνητρα και Μοντέλα Κύκλου Ζωής Επιθέσεων}\label{ch:3}
\chaptermeta{Ταξινόμηση δραστών · Το φάσμα των χάκερ · Cyber Kill Chain · MITRE ATT\&CK · Μοντέλο Διαμαντιού και D3FEND · Πληροφορίες απειλών και ανταλλαγή τους}{Επίπεδο: Εισαγωγικό–Μεσαίο επίπεδο · Σκέψη του αντιπάλου}{Χρόνος μελέτης: 6–8 ώρες}

Οι άμυνες αποκτούν νόημα μόνο σε σχέση με τους αντιπάλους που καλούνται να σταματήσουν. Μια μικρή επιχείρηση και ένας εθνικός πάροχος ενέργειας αντιμετωπίζουν πολύ διαφορετικούς αντιπάλους, με διαφορετικούς στόχους, πόρους και υπομονή. Για τον λόγο αυτό, το κεφάλαιο μετατοπίζει την οπτική από το σύστημα στον επιτιθέμενο.

Αρχικά ταξινομούμε τους δράστες απειλών με βάση το κίνητρο και τις ικανότητές τους και αποσαφηνίζουμε τις ηθικές και νομικές διακρίσεις πίσω από τη γνωστή ορολογία των «καπέλων». Στη συνέχεια μελετάμε τρία συμπληρωματικά μοντέλα που περιγράφουν την εξέλιξη μιας επίθεσης: το Cyber Kill Chain της Lockheed Martin, τη βάση γνώσης MITRE ATT\&CK και το Μοντέλο Διαμαντιού για την ανάλυση εισβολών. Το κεφάλαιο ολοκληρώνεται με τις πληροφορίες κυβερνοαπειλών (Cyber Threat Intelligence) —πώς παράγεται, αξιολογείται και ανταλλάσσεται υπεύθυνα η γνώση για τους αντιπάλους.

\begin{outcomesbox}{Μαθησιακά αποτελέσματα}
Μετά τη μελέτη του κεφαλαίου θα πρέπει να μπορείτε:
\begin{itemize}
  \item Να ταξινομείτε τους δράστες απειλών με βάση το κίνητρο, τις ικανότητες και την επιμονή τους και να τους συσχετίζετε με κατάλληλες άμυνες.
  \item Να διακρίνετε τη δραστηριότητα white-, grey- και black-hat με κριτήριο την εξουσιοδότηση και την πρόθεση.
  \item Να αντιστοιχίζετε μια επίθεση στις φάσεις του Cyber Kill Chain και στις τακτικές και τεχνικές του ATT\&CK.
  \item Να χρησιμοποιείτε το Μοντέλο Διαμαντιού για τη δομημένη ανάλυση μιας εισβολής.
  \item Να περιγράφετε τον κύκλο ζωής των πληροφοριών απειλών και να εφαρμόζετε το Πρωτόκολλο Φωτεινού Σηματοδότη (TLP) κατά την ανταλλαγή τους.
\end{itemize}
\end{outcomesbox}

\section{Ταξινόμηση των δραστών απειλών}\label{sec:3.1}
Οι δράστες απειλών ομαδοποιούνται χρήσιμα με βάση το τι επιδιώκουν και τι μπορούν να κάνουν. Οι \textbf{κρατικά υποστηριζόμενες ομάδες}, γνωστές και ως \emph{προηγμένες επίμονες απειλές} (Advanced Persistent Threats, APT), επιδιώκουν κατασκοπεία, δολιοφθορά ή επηρεασμό· διαθέτουν τους πόρους για εκμεταλλεύσεις μηδενικής ημέρας (zero-day) και για επιχειρήσεις μέσω της εφοδιαστικής αλυσίδας και μπορούν να παραμένουν σε ένα δίκτυο για μήνες ή χρόνια. Οι \textbf{κυβερνοεγκληματίες} έχουν οικονομικά κίνητρα· η οικονομία του λυτρισμικού ως υπηρεσίας (Ransomware-as-a-Service) έχει δώσει ακόμη και σε συνεργάτες μέτριας δεξιότητας πρόσβαση σε επαγγελματικά εργαλεία. Οι \textbf{χακτιβιστές} επιδιώκουν δημοσιότητα για έναν πολιτικό ή κοινωνικό σκοπό, συνήθως μέσω επιθέσεων άρνησης υπηρεσίας, αλλοίωσης ιστοτόπων ή διαρροών.

Οι \textbf{εσωτερικοί δράστες} χρειάζονται ιδιαίτερη προσοχή, διότι διαθέτουν ήδη νόμιμη πρόσβαση και, επομένως, ακυρώνουν τις παραδοχές που βασίζονται στην περίμετρο. Ένας \emph{κακόβουλος} εσωτερικός δράστης μπορεί να υποκλέψει δεδομένα ή να δολιοφθείρει συστήματα από απληστία ή μνησικακία, ενώ ένας \emph{αμελής} προκαλεί ζημιά ακούσια, για παράδειγμα μέσω λανθασμένης ρύθμισης ή επειδή ενέδωσε σε ένα μήνυμα ηλεκτρονικού ψαρέματος. Η κατανόηση των δραστών που αφορούν πραγματικά έναν οργανισμό τού επιτρέπει να επιλέγει μέτρα ανάλογα με τον πραγματικό κίνδυνο και όχι με τους πιο εντυπωσιακούς τίτλους ειδήσεων.

\begin{table}[htbp]
  \centering\small
  \caption{Προφίλ δραστών απειλών}
  \begin{tabularx}{\linewidth}{>{\raggedright\arraybackslash}X >{\raggedright\arraybackslash}X >{\raggedright\arraybackslash}X >{\raggedright\arraybackslash}X >{\raggedright\arraybackslash}X}
    \toprule
    \textbf{Δράστης} & \textbf{Κίνητρο} & \textbf{Ικανότητα} & \textbf{Τυπική διάρκεια παραμονής} & \textbf{Χαρακτηριστικές τεχνικές} \\
    \midrule
    Κρατική ομάδα APT & Κατασκοπεία, δολιοφθορά, επηρεασμός & Πολύ υψηλή & Μήνες έως χρόνια & Στοχευμένο ψάρεμα, χρήση νόμιμων εργαλείων του συστήματος, κρυφή διοίκηση και έλεγχος \\
    Κυβερνοεγκληματίας & Οικονομικό όφελος & Μέτρια έως υψηλή & Ημέρες έως εβδομάδες & Ψάρεμα, λυτρισμικό, κλοπή δεδομένων για εκβιασμό \\
    Χακτιβιστής & Ιδεολογία, δημοσιότητα & Χαμηλή έως μέτρια & Ώρες έως ημέρες & DDoS, αλλοίωση ιστοτόπων, δημοσιοποίηση προσωπικών στοιχείων \\
    Κακόβουλος εσωτερικός & Όφελος, μνησικακία & Μέτρια (εξουσιοδοτημένη πρόσβαση) & Μεταβλητή & Κατάχρηση προνομίων, διαρροή δεδομένων \\
    Αμελής εσωτερικός & Κανένα (σφάλμα) & Χαμηλή (ακούσια) & — & Λανθασμένες ρυθμίσεις, επαναχρησιμοποίηση κωδικών \\
    \bottomrule
  \end{tabularx}
\end{table}
\section{Το φάσμα των χάκερ: ηθική και νομιμότητα}\label{sec:3.2}
Η δημοφιλής ορολογία των «καπέλων» ταξινομεί την \textbf{εξουσιοδότηση και την πρόθεση}, όχι την τεχνική δεξιότητα. Ένας \textbf{white-hat} (ηθικός) χάκερ ελέγχει συστήματα μόνο με ρητή, γραπτή άδεια και εντός συμφωνημένου πεδίου, αναφέροντας τα ευρήματα ώστε να διορθωθούν. Ένας \textbf{black-hat} χάκερ ενεργεί χωρίς εξουσιοδότηση και με κακόβουλη ή ιδιοτελή πρόθεση. Ανάμεσά τους βρίσκεται ο \textbf{grey-hat}, ο οποίος μπορεί να ενεργεί χωρίς άδεια αλλά χωρίς σαφή κακοβουλία —για παράδειγμα, σαρώνει τον ιστότοπο μιας εταιρείας και στη συνέχεια την ενημερώνει για μια ευπάθεια.

Οι φοιτητές πρέπει να κατανοήσουν ότι οι καλές προθέσεις δεν δημιουργούν νομική εξουσιοδότηση. Στις περισσότερες έννομες τάξεις, συμπεριλαμβανομένων εκείνων που ενσωματώνουν την Οδηγία της ΕΕ για τις επιθέσεις κατά συστημάτων πληροφοριών, η πρόσβαση σε σύστημα χωρίς άδεια αποτελεί ποινικό αδίκημα ακόμη κι αν δεν προκληθεί ζημιά. Γι' αυτό η επαγγελματική δοκιμή διείσδυσης (Κεφάλαιο 10) ξεκινά πάντοτε με υπογεγραμμένο πεδίο εφαρμογής και κανόνες εμπλοκής, και γι' αυτό υπάρχουν τα προγράμματα συντονισμένης γνωστοποίησης ευπαθειών και τα προγράμματα επιβράβευσης (bug bounties): μετατρέπουν την περιέργεια του grey-hat σε εξουσιοδοτημένη και νομικά προστατευμένη έρευνα.

\section{Το Cyber Kill Chain της Lockheed Martin}\label{sec:3.3}
Το \textbf{Cyber Kill Chain}, που δημοσίευσε η Lockheed Martin το 2011, περιγράφει μια στοχευμένη εισβολή ως επτά διαδοχικές φάσεις: \textbf{αναγνώριση} (συλλογή πληροφοριών για τον στόχο), \textbf{οπλοποίηση} (σύζευξη μιας εκμετάλλευσης με ένα ωφέλιμο φορτίο), \textbf{παράδοση} (για παράδειγμα, ένα συνημμένο ηλεκτρονικού ταχυδρομείου), \textbf{εκμετάλλευση} (ενεργοποίηση της ευπάθειας), \textbf{εγκατάσταση} (εδραίωση μόνιμης παρουσίας), \textbf{διοίκηση και έλεγχος} (δημιουργία απομακρυσμένου καναλιού) και \textbf{ενέργειες επί των στόχων} (κλοπή δεδομένων, κρυπτογράφηση, δολιοφθορά).

Η βασική ιδέα του μοντέλου είναι αμυντική: επειδή ο επιτιθέμενος πρέπει να επιτύχει σε \emph{κάθε} φάση, ο αμυνόμενος αρκεί να σπάσει την αλυσίδα \emph{μία} φορά. Η έγκαιρη διακοπή είναι φθηνότερη —ο αποκλεισμός ενός κακόβουλου συνημμένου στη φάση της παράδοσης αποτρέπει μια δαπανηρή απόκριση σε περιστατικό αργότερα. Το μοντέλο έχει όμως περιορισμούς. Σχεδιάστηκε με βάση εισβολές με κακόβουλο λογισμικό που προέρχονται από το εξωτερικό και περιγράφει λιγότερο ικανοποιητικά τις εσωτερικές απειλές, την κατάχρηση διαπιστευτηρίων και τις επιθέσεις στο νέφος. Για τον λόγο αυτό συνδυάζεται συνήθως με το λεπτομερέστερο πλαίσιο ATT\&CK.

\section{MITRE ATT\&CK: τακτικές, τεχνικές και διαδικασίες}\label{sec:3.4}
Το \textbf{MITRE ATT\&CK} είναι μια δημόσια διαθέσιμη βάση γνώσης για τη συμπεριφορά των αντιπάλων, η οποία έχει δημιουργηθεί από παρατηρήσεις πραγματικών εισβολών. Οργανώνεται ως πίνακας. Οι στήλες είναι οι \textbf{τακτικές} —οι βραχυπρόθεσμοι στόχοι του επιτιθέμενου, όπως η \emph{Αρχική Πρόσβαση}, η \emph{Εκτέλεση}, η \emph{Μόνιμη Παρουσία}, η \emph{Κλιμάκωση Προνομίων}, η \emph{Αποφυγή Άμυνας}, η \emph{Πρόσβαση σε Διαπιστευτήρια}, η \emph{Πλευρική Κίνηση} και η \emph{Διαρροή Δεδομένων}. Κάθε στήλη περιέχει \textbf{τεχνικές} που περιγράφουν \emph{πώς} επιτυγχάνεται ο στόχος· για παράδειγμα, η T1059.001 δηλώνει εκτέλεση εντολών μέσω PowerShell. Οι \textbf{διαδικασίες} είναι οι συγκεκριμένοι τρόποι με τους οποίους μια ομάδα έχει χρησιμοποιήσει μια τεχνική.

Το ATT\&CK παρέχει στους αμυνόμενους μια κοινή γλώσσα. Μια ομάδα ασφάλειας μπορεί να αντιστοιχίσει τους κανόνες ανίχνευσής της σε τεχνικές και να διαπιστώσει αμέσως πού υπάρχουν κενά κάλυψης· μια αναφορά πληροφοριών απειλών μπορεί να δηλώνει ποιες τεχνικές χρησιμοποιεί μια ομάδα· και μια κόκκινη ομάδα (red team) μπορεί να προσομοιώσει αυτές τις τεχνικές για να ελέγξει αν οι άμυνες λειτουργούν πράγματι. Η έμφαση στη \emph{συμπεριφορά} και όχι σε εύκολα μεταβαλλόμενες ενδείξεις, όπως οι τιμές κατακερματισμού αρχείων, είναι αυτό που καθιστά το πλαίσιο ανθεκτικό στον χρόνο.

\section{Το Μοντέλο Διαμαντιού και το MITRE D3FEND}\label{sec:3.5}
Το \textbf{Μοντέλο Διαμαντιού για την Ανάλυση Εισβολών} (Diamond Model) αναπαριστά κάθε κακόβουλο γεγονός μέσω τεσσάρων συνδεδεμένων στοιχείων: του \textbf{αντιπάλου}, της \textbf{δυνατότητας} (εργαλεία και κακόβουλο λογισμικό), της \textbf{υποδομής} (ονόματα τομέα, διακομιστές, διευθύνσεις IP) και του \textbf{θύματος}. Η δύναμη του μοντέλου βρίσκεται στη \emph{μετάβαση} (pivoting) από στοιχείο σε στοιχείο: μόλις ο αναλυτής γνωρίζει μία κορυφή, οι συνδέσεις υποδεικνύουν πού να αναζητήσει στη συνέχεια. Ένα δείγμα κακόβουλου λογισμικού (δυνατότητα) μπορεί να αποκαλύψει έναν τομέα διοίκησης και ελέγχου (υποδομή), του οποίου τα στοιχεία καταχώρισης μπορεί να οδηγούν σε άλλα θύματα και, τελικά, σε ένα προφίλ αντιπάλου.

Ενώ το ATT\&CK καταγράφει την επιθετική συμπεριφορά, το \textbf{MITRE D3FEND} καταγράφει αμυντικά αντίμετρα —όπως η \emph{ανάλυση δημιουργίας διεργασιών} ή η \emph{θωράκιση διαπιστευτηρίων}— και τα συνδέει με τις επιθετικές τεχνικές που αντιμετωπίζουν. Σε συνδυασμό, το Kill Chain απαντά στο \emph{πότε} να παρέμβουμε, το ATT\&CK στο \emph{τι} κάνει ο αντίπαλος, το Μοντέλο Διαμαντιού στο \emph{ποιος και με ποια μέσα} και το D3FEND στο \emph{ποιο μέτρο} αντιμετωπίζει κάθε συμπεριφορά.

\section{Πληροφορίες κυβερνοαπειλών: κύκλος ζωής, είδη και υπεύθυνη ανταλλαγή}\label{sec:3.6}
Οι \textbf{πληροφορίες κυβερνοαπειλών} (Cyber Threat Intelligence, CTI) είναι πληροφορίες για τους αντιπάλους που έχουν συλλεχθεί, αναλυθεί και τοποθετηθεί σε πλαίσιο, ώστε να υποστηρίζουν τη λήψη αποφάσεων. Παράγονται μέσω ενός \textbf{κύκλου ζωής πληροφοριών}: \emph{καθοδήγηση} (προσδιορισμός του τι χρειάζεται να γνωρίζουν όσοι λαμβάνουν αποφάσεις), \emph{συλλογή}, \emph{επεξεργασία}, \emph{ανάλυση}, \emph{διάχυση} και \emph{ανατροφοδότηση}. Οι πληροφορίες διακρίνονται συχνά σε τρία επίπεδα. Οι \textbf{στρατηγικές} ενημερώνουν τη διοίκηση για τάσεις και κινδύνους· οι \textbf{επιχειρησιακές} περιγράφουν εκστρατείες και δράστες· οι \textbf{τακτικές} παρέχουν τεχνικές λεπτομέρειες, όπως τεχνικές επίθεσης και ενδείξεις παραβίασης.

Οι πληροφορίες αποκτούν αξία όταν ανταλλάσσονται, όμως η ανταλλαγή πρέπει να σέβεται τις επιθυμίες της πηγής. Το \textbf{Πρωτόκολλο Φωτεινού Σηματοδότη} (Traffic Light Protocol, TLP 2.0) χαρακτηρίζει τις πληροφορίες ως \emph{TLP:RED} (μόνο ονομαστικοί παραλήπτες), \emph{TLP:AMBER} (ο οργανισμός του παραλήπτη, βάσει ανάγκης γνώσης), \emph{TLP:GREEN} (η ευρύτερη κοινότητα) ή \emph{TLP:CLEAR} (δημόσιες). Η μηχανικά αναγνώσιμη ανταλλαγή βασίζεται στο \textbf{STIX}, μια τυποποιημένη γλώσσα περιγραφής πληροφοριών απειλών, και στο \textbf{TAXII}, ένα πρωτόκολλο μεταφοράς τους. Κλαδικές κοινότητες, όπως τα Κέντρα Ανταλλαγής και Ανάλυσης Πληροφοριών (ISAC), επιτρέπουν σε οργανισμούς που αντιμετωπίζουν παρόμοιες απειλές να προειδοποιούν γρήγορα ο ένας τον άλλον.

\section*{Βασικοί όροι}
\begin{description}[style=nextline,leftmargin=1.2em]
  \item[APT] Προηγμένη επίμονη απειλή: δράστης με σημαντικούς πόρους που διατηρεί μακροχρόνια κρυφή πρόσβαση.
  \item[Kill Chain] Μοντέλο επτά φάσεων μιας στοχευμένης εισβολής, από την αναγνώριση έως τις ενέργειες επί των στόχων.
  \item[TTP] Τακτικές, τεχνικές και διαδικασίες: η συμπεριφορική περιγραφή του τρόπου δράσης ενός αντιπάλου.
  \item[Μοντέλο Διαμαντιού] Αναλυτικό μοντέλο που συνδέει αντίπαλο, δυνατότητα, υποδομή και θύμα.
  \item[TLP] Πρωτόκολλο Φωτεινού Σηματοδότη: χαρακτηρισμοί που ορίζουν πόσο ευρέως μπορεί να αναδιανεμηθεί μια πληροφορία.
  \item[STIX / TAXII] Τυποποιημένη γλώσσα (STIX) και πρωτόκολλο μεταφοράς (TAXII) για την ανταλλαγή πληροφοριών απειλών.
\end{description}

\section*{Σύνοψη κεφαλαίου}
\begin{itemize}
  \item Οι δράστες απειλών διαφέρουν ως προς το κίνητρο, τις ικανότητες και την επιμονή· οι άμυνες πρέπει να είναι ανάλογες με τους δράστες που πράγματι αφορούν τον οργανισμό.
  \item Η ορολογία των «καπέλων» αφορά την εξουσιοδότηση και την πρόθεση· η μη εξουσιοδοτημένη πρόσβαση είναι παράνομη ανεξάρτητα από το κίνητρο.
  \item Το Kill Chain δείχνει ότι η διακοπή οποιασδήποτε φάσης σταματά την εισβολή και ότι η έγκαιρη διακοπή κοστίζει λιγότερο.
  \item Το ATT\&CK περιγράφει λεπτομερώς τη συμπεριφορά των αντιπάλων και παρέχει κοινή γλώσσα για την ανίχνευση, τις πληροφορίες απειλών και τις δοκιμές.
  \item Το Μοντέλο Διαμαντιού υποστηρίζει τη μετάβαση μεταξύ στοιχείων κατά την ανάλυση, ενώ οι πληροφορίες πρέπει να ανταλλάσσονται με σαφείς κανόνες, όπως το TLP.
\end{itemize}

\section*{Ερωτήσεις ανασκόπησης}
\begin{enumerate}
  \item Γιατί οι εσωτερικές απειλές αντιμετωπίζονται δύσκολα με μέτρα που βασίζονται στην περίμετρο; Προτείνετε δύο μέτρα αποτελεσματικά απέναντί τους.
  \item Αντιστοιχίστε μια τυπική επίθεση που ξεκινά με ηλεκτρονικό ψάρεμα και καταλήγει σε λυτρισμικό στις επτά φάσεις του Cyber Kill Chain.
  \item Εξηγήστε γιατί οι ανιχνεύσεις που βασίζονται σε τεχνικές του ATT\&CK είναι πιο ανθεκτικές από εκείνες που βασίζονται σε τιμές κατακερματισμού αρχείων.
  \item Ένας συνεργάτης σάς αποστέλλει αναφορά με χαρακτηρισμό TLP:AMBER. Μπορείτε να την προωθήσετε σε έναν προμηθευτή; Τεκμηριώστε την απάντησή σας.
\end{enumerate}
